Write \(\tau_m=\tau_{c^\dagger}(m)\).  All divisions by \(\pi_r^{n,M}\) below are restricted to \(\pi_r^{n,M}>0\); the rule on a zero-mass allocation orbit is immaterial.  Since \(\kappa_{r,x}^{n,M}\) is a probability distribution on \(\Gamma_M\subset[-1,1]\), its barycenter \(\delta_{n,M}(r,x)\) is an admissible deterministic action.  For every \(m,r,x\), the conditional variance identity gives
\[
 \sum_g\kappa_{r,x}^{n,M}(g)(g-\tau_m)^2
 =\{\delta_{n,M}(r,x)-\tau_m\}^2
 +\sum_g\kappa_{r,x}^{n,M}(g)
 \{g-\delta_{n,M}(r,x)\}^2.
\]
The last term is nonnegative.  Multiplication by \(\pi_r^{n,M}P_m(x\mid r)\), followed by summation over \((r,x)\), proves
\[
 \rho_n^\dagger\leq U_{n,M}\leq\Lambda_{n,M}^\dagger.
\]
The first inequality uses the exact orbit-game equality: the barycentered orbit rule is an admissible labeled procedure obtained by uniform assignment conditional on the selected allocation count.

We next derive the lower certificate from the primitive likelihood.  Fix a prior \(\nu\) and an allocation orbit \(r\).  Put
\[
 Q_\nu=\sum_m\nu_m\tau_m^2,
 \qquad D_{r,x}=D_{r,x}(\nu),
 \qquad N_{r,x}=N_{r,x}(\nu).
\]
If \(D_{r,x}=0\), nonnegativity of every summand implies \(N_{r,x}=0\), so the action at \(x\) has no effect on Bayes risk.  If \(D_{r,x}>0\), completing the square yields
\[
 \begin{aligned}
 &\sum_m\nu_m\sum_xP_m(x\mid r)\{\delta(r,x)-\tau_m\}^2\\
 &\quad=Q_\nu-
 \sum_{x:D_{r,x}>0}\frac{N_{r,x}^2}{D_{r,x}}
 +\sum_{x:D_{r,x}>0}D_{r,x}
 \left\{\delta(r,x)-\frac{N_{r,x}}{D_{r,x}}\right\}^2.
 \end{aligned}
\]
Thus the displayed posterior mean minimizes Bayes risk at fixed \(r\), and the minimum is the bracket defining \(B_n(\nu)\).  Moreover, for an arbitrary design mixture \(\pi\), minimization over all orbit rules gives
\[
 \inf_\delta\sum_m\nu_mR_m(\pi,\delta)
 =\sum_r\pi_rb_r(\nu)\geq\min_rb_r(\nu)=B_n(\nu).
\]
Randomization over allocation orbits therefore cannot lower Bayes risk below \(B_n(\nu)\).  Distribute mass \(\nu_m\) uniformly over the labeled schedules in orbit \(m\).  This is an invariant prior on fixed schedules.  The symmetrization theorem applies to every labeled design and estimator, so maximum labeled-schedule risk dominates this prior Bayes risk.  Hence
\[
 B_n(\nu)\leq\rho_n^\dagger.
\]
No positivity condition is needed: zero allocation probabilities and zero posterior denominators have already been handled explicitly.

It remains to compare the grid and continuous Bayes actions.  For fixed \(\nu,r,x\) with \(D_{r,x}>0\), completing the same square but restricting the action to \(\Gamma_M\) adds
\[
 D_{r,x}\operatorname{dist}\!\left(
 \frac{N_{r,x}}{D_{r,x}},\Gamma_M\right)^2.
\]
Let \(B_{n,M}^{\mathrm{grid}}(\nu)\) be the minimum over \(r\) of the sum of this term and \(b_r(\nu)\).  In the dual of the finite rational program, the multipliers on the response-count risk inequalities are nonnegative and sum to one by the stationarity equation for the epigraph variable.  Eliminating the dual variables for the occupancy constraints gives exactly
\[
 \Lambda_{n,M}^\dagger
 =\max_{\nu\in\Delta(\mathcal M_{n,3})}
 B_{n,M}^{\mathrm{grid}}(\nu).
\]
The rational primal--dual attainment used here is supplied by the grid-LP certificate proposition.  For the certified maximizer \(\nu_{n,M}\), therefore,
\[
 \Lambda_{n,M}^\dagger=B_{n,M}^{\mathrm{grid}}(\nu_{n,M}).
\]
The action interval is \([-1,1]\), the grid mesh is \(1/M\), and hence every posterior mean is within \(1/(2M)\) of \(\Gamma_M\).  Also
\[
 \sum_xD_{r,x}(\nu)=\sum_m\nu_m\sum_xP_m(x\mid r)=1.
\]
Consequently
\[
 B_n(\nu)\leq B_{n,M}^{\mathrm{grid}}(\nu)
 \leq B_n(\nu)+\frac{1}{4M^2},
\]
which completes the finite-sample sandwich.

For the asymptotic statement, substitute \(M=M_n\) and multiply the sandwich by \(a_n\).  The two certificate improvements bracket \(a_nd_n(c^\dagger)\), and their difference is at most \(a_n/(4M_n^2)\to0\).  If the bracketed middle sequence has a finite limit, both endpoints have the same limit by squeezing.

Finally, evaluability is finite and explicit.  There are \(\binom{n+7}{7}\) response-count constraints, \(\binom{n+2}{2}\) allocation orbits, and
\[
 \sum_{r_1+r_2+r_3=n}(r_1+1)(r_2+1)(r_3+1)=\binom{n+5}{5}
\]
allocation--observation orbits.  Thus the grid certificate is obtained by one rational linear program with
\[
 1+\binom{n+2}{2}+(2M+1)\binom{n+5}{5}
\]
primal variables, followed by finite rational sums and maxima for \(U_{n,M}\) and \(B_n(\nu_{n,M})\).